Seit der Veröffentlichung der moderne Portfoliotheorie in 1952, wofür Harry Markowitz einen Nobelpreis für Wirtschaft erhalten hat, hat sich die Vermögensverwaltung grundlegend geändert. Der mathematische Beweis dafür, dass Portfoliorisiko durch die Hinzunahme nicht-perfekt-korrelierter Assets ins Portfolio reduziert werden kann, hat die moderne Portfoliotheorie von Markowitz einen Platz als systematischer Ansatz zur Lösung des Assetallokationsproblem gegeben. Allerdings leiden die Lösungen des mathematischen Optimierungsproblem zur Berechnung von Markowitz optimalen Portfolien an niedrige Streuung, schlechte out-of-sample Performance und Instabilität bzgl. Änderungen der Inputparametern. In dieser Arbeit geben wir einen Überblick des Assetallokationsproblems, definieren das mathematische Optimierungsproblem zur Berechnung von Markowitz optimalen Portfolien und analysieren die Grunde und Quellen der Instabilität. Au{\ss}erdem, zeigen wir drei Methoden, die die Instabilität reduzieren sollen: zwei beschäftigen sich mit der Kovarianzmatrix, die als Parameter ins Optimierungsproblem reinfliesst, und die dritte Methode ist ein Alternativansatz zu den Assetallokationsproblem. Zuletzt vergleichen wir die Performance bzgl. Streuung, out-of-sample risiko-adjustierter Rendite und zeitlicher Stabilität Markowitz optimaler Portfolio mit der der anderen drei Methoden während des Jahres 2020, als die COVID-19 Pandemie ausbrach.